\chapter{ Introduction }
\label{sec:intro}
%Things the intro needs to have- ideally in this order
%General State of the world
%problem? What is wrong with the general state of the world?
%optional but useful- Insight
%Approach- what are you doing to address this problem AND why is this reasonable?
%Relation to prior work- how does your approach differ from/ build on what has already been done- this estabilshes both novelty and relevance.
%Evaluation/ results- what did you do to demonstrate your approach shows promise (in brief)
%Implications- what are the broader implications of the results
%explicit statement of contributions

%  The general state of the world is that each embedded audio target has its own siloed software ecosystem, maintained by the hardware manufacturer.

% Problems with this:
% \begin{itemize}
% \item The quality of each indiviudal algorithm is diminished by the fact that the hardware mfg must crank out all possible audio processing that end users expect.
% \item End users often cannot have a uniformity of sound across desktop and embedded audio contexts
% \item Shifts in popular audio DSP endebt the mfg to continually push firmware updates with new processing
% \end{itemize}

% My insight is to use a target-agnostic IR as the distributable for audio DSP code, so that the same DSP can be used across any target that has a runtime for it. This is not a new concept in application software generally, but has not been applied to this area.

%  My approach is to benchmark and compare three runtimes for the WASM on an embedded target. The embedded target (Daisy Seed) is a popular choice for developers in embedded audio, and representative of the type of processor found in many low-cost digital audio effect devices. The algorithms chosen represent major pillars of popular audio effect DSP: white-box modeling, black-box modeling, and phsyical modeling. Notably, NAM represents a "canary in the coalmine" moment for embedded audio- when the open standard was released, hardware mfgs scrambled to add support to their platforms.

% My relationship to prioir works is essentially fusing research from two fields. In the audio DSP world, tools like iPlug2 and FAUST have been created to target multiple platforms from one codebase. In IoT research, WASM is being explored as a "write once, run anywhere" for sandboxed application software.

% My evaluation is that runtime interpreters for audio DSP encapsulated as WASM is not viable- but translation of the WASM to native code at flash-time offers acceptable performance.

% The implication is that a well crafted WASM ABI for headless audio plugins can be targeted from many programming languages, developed and tested in a client-side web IDE, and distributed for use in multiple environments.

% My explicit contribution is the explore avenues of translation and the benchmarking of their relative performance, as well as insights into potential next steps.

The landscape of audio software in 2026 is diverse. Deployments include standalone applications spanning various operating systems and computer architectures, myriad plugin formats for digital audio workstations (DAWs), game engine plugins, web audio, embedded systems, and mobile apps, to name a few. In many cases, interoperability between these environments is desirable from both end user and developer perspectives.

For users, interoperability promotes sonic consistency. For example, an electric instrumentalist may want the same effect processing they use in a DAW to also be available in a live performance context, ideally on hardware that is tailored to that environment. For developers, deployment to multiple targets from a single codebase and build system endows greater ability to reach more end users with their software.

 % Within the audio developer community, many solutions have arisen to facilitate deployment to multiple targets from a single codebase, although few are tailored to embedded systems and the vast majority rely on conditional builds for each target.


% All figures should be centered on the column (or page, if the figure spans both
% columns). Figure captions (in italic) should follow each figure and have the
% format given in Figure~\ref{fft_plot}. Vectorial figures are preferred (e.g.,
% Postscript, PDF, etc.). Also, in order to provide a better readability, figure
% text font size should be at least identical to footnote font size. If bitmap
% figures are used, please make sure that the resolution is enough for print
% quality. Figure~\ref{ftt_plot2} illustrates an example of a figure spanning two
% columns.

The goal of this project is to investigate means of targeting myriad deployments from one (non-source) distributable provided by audio developers, a \textit{universal distributable}. Such a deliverable can confer many advantages, but also carries with it some inherent downsides. Advantages include simplified build systems, easier porting of algorithms to new platforms, and greater interoperability between heterogeneous deployments. Downsides include performance degradation from the abstraction layer, increased complexity of debugging, and easier de-compilation. In this project I have focused on deployment of such a universal distributable in embedded systems, as these resource-contained environments provide the greatest challenge in portability. If deployment to embedded systems proves feasible, then all other targets may follow suit with relative ease.

\section { Embedded Audio Systems }

Embedded systems are computer systems specialized for a particular task\cite{catsoulis2005designing}. Examples that perform audio processing surround us: smart home devices, digital video cameras, and wireless earphones. Common form factors for musical contexts include sampling synthesizers, stompbox effect processors, and digital mixers.

Traditionally, these systems are developed with tightly controlled, manufacturer-specific software ecosystems\cite{10.1145/266021.266373}. Each hardware platform provides its own developer tooling, firmware architecture, and collection of DSP modules. While this tight coupling between software and hardware provides advantages in performance, it also produces ecosystem siloing. Algorithms are not easily portable across devices, and innovation is constrained by the resources and priorities of individual manufacturers.

This ecosystem siloing introduces several practical limitations. Most notable among them:

\begin{itemize}
    \item Individual algorithm quality may suffer because hardware vendors are responsible for implementing and maintaining the broad range of processing that users expect, often redundantly \textit{re-}implementing the popular modules offered by similar products.

    \item Users cannot easily maintain sonic consistency between desktop environments and embedded devices, as the robust world of third-party plugin software for desktop audio is not available in embedded contexts.

    \item Shifts in popular audio processing techniques create ongoing maintenance burdens. When new paradigms emerge, manufacturers must issue firmware updates or even redesign hardware to remain competitive.
\end{itemize}

Given the resource-constrained nature of embedded systems, the goal of using a unified distributable to address the above listed issues is lofty- but not unprecedented.

% The recent advent of Neural Amp Modeler[?] (NAM), in particular, illustrates a critical inflection point for embedded audio systems: upon release of its open standard, hardware manufacturers scrambled to integrate support, highlighting the limitations of closed, siloed architectures.

% Need sources for above claims.

%  ^ I recognize that the above sentence on NAM is gratuitous... but I'm inclined to mentiond it somewhere.

\section { WebAssembly as an Intermediate Representation }

% To remedy the issues described above, I propose distributing audio DSP modules in a target-agnostic intermediate representation (IR), allowing the same encapsulated DSP to be deployed on any platform that supports a corresponding runtime\footnote{The term "runtime" is used liberally here, and will be explained further in section 2.2}.

% Although the concept of portable intermediate representation is well established in general application software, research into its use for embedded audio DSP is yet to gain major traction.

In computing generally, intermediate representation (IR) often takes one of two forms. Some forms, such as LLVM IR\cite{1281665}, are used internally by compilers as a translation stage in the process of transforming high-level source code into architecture-specific machine code prior to program execution. Another form of IR is \textit{bytecode}, which is also translated from high-level source code but is intended to run inside a \textit{virtual machine}\cite{lindholm2013java}, a software program that translates IR to native machine code at runtime.

% \footnote{The term \textit{bytecode} comes from varieties of IR that encode operations as one-byte words. Some forms of bytecode however, including WASM, use binary (bit) encoding and consequently are sometimes referred to as \textit{bitcode}. In this project I use the term bytecode as a catch-all for any IR intended for use with an interpreter.}

 Recent Internet of Things (IoT) research\cite{fi15080275}\cite{10756513}\cite{IOTW753} has explored WebAssembly (WASM) as an IR of both forms described above, with the general goal of WASM as a "write once, run anywhere" substrate for self-contained applications. A major advantage of WASM over other IRs is that the "write once" step can occur in many different programming languages, as the standard was developed from the onset to be language-agnostic\cite{og_wasm}.

% When used one way, it's interpretation, another way, AOT,

Prior research in digital audio has proposed WASM as a universal distributable for audio DSP with a framework entitled \textit{Audio Anywhere}\cite{gaster2020audio}, although development of this project has ceased. In their work, the authors explored interpreter-based execution and ahead-of-time (AOT) translation strategies for their runtimes targeting embedded, but limited their scope to WASM modules compiled from a single input source language. Building on their research, I have investigated alternative AOT translation strategies for WASM modules, enabled additional high-level source code languages (C/C++) as input, benchmarked a wider range of audio processing algorithms, and developed web-based tooling for development and distribution.

 % , and focus exclusively on embedded contexts.

% A particular advantage of WASM is that is it LLVM-based, so many programming langauges can target its runtimes.

% To assess the viability of my approach, I have measured the performance of three WASM runtime strategies and compared against native execution on an embedded target representative of low-cost digital audio hardware: the Daisy Seed microcontroller platform. The selected DSP algorithms span major paradigms in contemporary audio effects design, including white-box, black-box, and physical modeling.

Inline with performance measurements presented in \cite{gaster2020audio}, my benchmark results indicate that runtime interpretation of WASM is not a viable strategy for real-time embedded audio workloads. However, runtimes leveraging AOT translation of WASM modules yield acceptable performance under certain conditions, notably when a program's memory footprint is known at compile time. Performance may be further improved by better leveraging WASM's interface for importing functions from its host, greater optimization of compiler toolchains, and increased runtime specialization.

% Notably, one of the AOT methods I utilize does not depend on local development tools, as I have embedded it in a client-side web application. <- put this somewhere else

My findings suggest that with further research, a carefully specified WASM ABI for audio processing and parameter management could enable language-agnostic development of encapsulated audio DSP coupled with deployment across embedded devices, desktop applications, and web platforms from a universal distributable. In that aim, the primary contributions of this work are the exploration of multiple execution strategies for WASM-based audio DSP modules, empirical benchmarking of their relative performance on representative embedded hardware, and creation of web-based developer tooling to facilitate their authoring and deployment.
